\documentclass[12pt,reqno]{amsart}
\usepackage{math-review}
% Project-wide mathematical notation. No claims are certified here.

\title{Strong Compactness and Admissible Limits -- Minimal}
\author{Known-mathematics exposition benchmark}
\date{}
\begin{document}
\maketitle
\section{Main Results}
The question starts with an existing solution sequence. The task is to
justify its limit, not to prove a priori that the sequence or estimates exist.

For an integer $d\geq1$, write $\rho$ for flat distance on
$\mathbb T^d=\mathbb R^d/\mathbb Z^d$ and set, for $0<\gamma<1$,
\[
 [v]_\gamma=\sup_{x\ne y}\frac{|v(x)-v(y)|}{\rho(x,y)^\gamma},\qquad
 \|u\|_{C^{2,\gamma}}
 =\max_{|\nu|\leq2}\|D^\nu u\|_\infty
   +\max_{|\nu|=2}[D^\nu u]_\gamma.
\]
Here $D^\nu$ denotes a periodic coordinate derivative of multi-index $\nu$.
Only the fixed flat structure is used below.

Matrix inequalities below mean inequalities of quadratic forms.
\begin{theorem}[Conditional sequential closedness]\label{thm:main}
Fix integers $d,r\geq1$ and $0<\beta<\alpha<1$.
On the flat torus $\mathbb T^d=\mathbb R^d/\mathbb Z^d$, let
\[
 A:C^{2,\beta}(\mathbb T^d)\longrightarrow
 C^0(\mathbb T^d;\operatorname{Sym}_r(\mathbb R))
\]
be continuous, where the target has the uniform operator norm, and set
\[
 \mathcal U=\{u\in C^{2,\beta}(\mathbb T^d):
       A(u)(x)>0\text{ for every }x\in\mathbb T^d\}.
\]
Let $\mathcal F:[0,1]\times\mathcal U\to C^0(\mathbb T^d)$ be continuous.
Suppose $t_j\in[0,1]$, $t_j\to t_\infty$, and
$u_j\in C^{2,\alpha}(\mathbb T^d)\cap\mathcal U$ satisfy, for constants
$C<\infty$ and $\delta>0$ independent of $j$,
\[
 \mathcal F(t_j,u_j)=0,\qquad
 \|u_j\|_{C^{2,\alpha}}\leq C,\qquad
 A(u_j)(x)\geq\delta I_r\quad\text{for every }j,x.
\]
Then a subsequence converges in $C^{2,\beta}$ to a function
$u_\infty\in C^{2,\alpha}\cap\mathcal U$, with
\[
 A(u_\infty)(x)\geq\delta I_r\quad\text{for every }x,
 \qquad \mathcal F(t_\infty,u_\infty)=0.
\]
The norms are the fixed periodic derivative/H\"older norms defined immediately
before this statement. This is a conditional sequential closedness assertion,
not an existence theorem or a smoothness assertion.
\end{theorem}
The three hypotheses doing different jobs are a strong compactness bound,
a uniform distance from the matrix-cone boundary, and continuity of the equation.

\section{Key Mathematical Ingredients}
The first input specifies what the estimate actually buys; a generic
reference to compactness would not specify a topology for the nonlinear limit.
\begin{samepage}
\begin{lemma}\label{lem:compactness}
On the flat torus $\mathbb T^d$ with fixed periodic derivative norms,
let $0<\beta<\alpha<1$. If a sequence of real functions satisfies
\[
 u_j\in C^{2,\alpha}(\mathbb T^d),\qquad
 \sup_j\|u_j\|_{C^{2,\alpha}}<\infty,
\]
then some subsequence converges strongly in $C^{2,\beta}(\mathbb T^d)$
to a function in $C^{2,\alpha}(\mathbb T^d)$.
The assertion does not claim convergence in the stronger $C^{2,\alpha}$ norm.
\end{lemma}
\end{samepage}
Here this classical compactness input is used, not reproved. Its derivative
extraction and interpolation proof is developed in the companion lecture.
The second input keeps the limit within the domain of the equation.
\begin{lemma}\label{lem:margin}
Let $H_j,H\in C^0(\mathbb T^d;\operatorname{Sym}_r(\mathbb R))$,
let $\delta>0$ be fixed, and assume $H_j\to H$ uniformly in operator norm.
If $H_j(x)\geq\delta I_r$ for every $j,x$, then $H(x)\geq\delta I_r$
for every $x$. For positive definite $H_j(x)$ the same uniform lower bound
is equivalent to
\[
 \|H_j(x)^{-1}\|_{\mathrm{op}}\leq\delta^{-1}
 \quad\text{for every }j,x.
\]
\end{lemma}
The first implication follows by testing a fixed unit vector and passing
to the uniform limit. The inverse formulation uses orthogonal diagonalization.

\section{Connecting the Ingredients}
Apply Lemma~\ref{lem:compactness} and relabel the resulting subsequence:
\[
 u_j\longrightarrow u_\infty\quad\text{in }C^{2,\beta}.
\]
Continuity of $A$ converts this into uniform convergence of matrices.
For every $x$ and every unit vector $v\in\mathbb R^r$,
\[
 v^{\mathsf T}A(u_\infty)(x)v
 =\lim_jv^{\mathsf T}A(u_j)(x)v\geq\delta.
\]
Consequently $u_\infty\in\mathcal U$. Only now may continuity on
$[0,1]\times\mathcal U$ be invoked to obtain
\[
 \mathcal F(t_\infty,u_\infty)
 =\lim_j\mathcal F(t_j,u_j)=0\quad\text{in }C^0.
\]
The uniform margin is indispensable to this route. With constant functions
$u_j=t_j=1/j$, take $A(u)=uI_r$ and $\mathcal F(t,u)=u-t$ on positive functions.
The functions are uniformly H\"older-bounded solutions, but their limit is zero
and is not admissible. Separate strict positivity is not a uniform lower bound.

\nocite{S-guide}
\bibliographystyle{plain}
\bibliography{../generated/references}
\end{document}
